% Part: proof-theory
% Chapter: cut-elimination
% Section: interpolation

\documentclass[../../../include/open-logic-section]{subfiles}

\begin{document}

\olfileid{pt}{cut}{itp}

\olsection{మహేరా లెమ్మా, క్రెయిగ్ అంతర్వర్తన ప్రమేయం}

విలియం క్రెయిగ్ ప్రతిపాదించిన అంతర్వర్తన ప్రమేయం ఇలా చెబుతుంది: $!A
\Entails !B$ అయితే $!A \Entails !C$, $!C \Entails !B$ అయ్యే \tetoken{ఒక వాక్యం}{!!a{sentence}}~$!C$ (\emph{మధ్యవర్తి వాక్యం}) ఉంది. ముఖ్యంగా,~$!C$లో కనిపించే \tetoken{స్థిరసంకేతాలు}{!!{constant}s}, \tetoken{విధేయ సంకేతాలు}{!!{predicate}s} అన్నీ~$!A$,~$!B$ రెండింటిలోనూ కనిపించేలా మధ్యవర్తి~$!C$ను ఎంచుకోవచ్చు. (ఇక్కడ \tetoken{ప్రమేయ సంకేతాలు}{!!{function}s} లేవని అనుకుంటాం.) ఈ పరిమితి లేకపోతే $!A$,~$!B$లనే మధ్యవర్తులుగా సులభంగా తీసుకోవచ్చు.

అంతర్వర్తన ప్రమేయం సంప్రదాయ మొదటిస్థాయి తర్కానికి వర్తిస్తుంది; దాన్ని కనుగొన్న “మొదటిస్థాయి తర్కపు చివరి ముఖ్య లక్షణం” అని కూడా పిలిచారు. నమూనా సిద్ధాంతం, స్వయంచాలక తర్కంలో దానికి ముఖ్య అన్వయాలున్నాయి. ఒక సిద్ధాంతాన్ని ఏ ప్రాథమిక భావాలతో స్వీకృతీకరించవచ్చో, ఏవాటితో చేయలేమో అధ్యయనం చేయడానికి వాడగలగడం వల్ల విజ్ఞాన తత్వశాస్త్రం దాని అసలు ప్రేరణ, అన్వయ రంగం. అంతేగాక, ఒక సిద్ధాంతం ఏ భావాన్నైనా పరోక్షంగా నిర్వచించగలిగితే ప్రత్యక్షంగానూ నిర్వచించగలదని చెప్పే బెత్ నిర్వచనీయత ప్రమేయం కూడా దీని నుంచి వస్తుంది.

గణిత తర్కంలోని అనేక ఫలితాలలాగే అంతర్వర్తన ప్రమేయాన్ని నమూనా సిద్ధాంత, నిరూపణ సిద్ధాంత పద్ధతులు రెండింటితోనూ నిరూపించవచ్చు. నిరూపణ సిద్ధాంత పద్ధతుల ప్రయోజనం ఏమిటంటే, అవి మధ్యవర్తి ఉనికిని చూపడమే కాక, $!A \Entails !B$కు గల \tetoken{ఒక నిరూపణ}{!!a{proof}} నుంచి మధ్యవర్తిని లెక్కించే అల్గోరిథాన్ని ఇస్తాయి; ఉదాహరణకు $!A \Sequent !B$కు~\Log{G3c}లో ఉన్న \tetoken{ఒక నిరూపణ}{!!a{proof}} నుంచి.

$\Lang L(!A)$ అనేది~$!A$లోని \tetoken{స్థిరసంకేతాలు}{!!{constant}s}, \tetoken{విధేయ సంకేతాల}{!!{predicate}s} సమితి; $\Lang L(\Gamma) = \bigcup \Setabs{\Lang L(!A)}{!A \in
\Gamma}$ అనుకుందాం. ఈ విభాగమంతటా \tetoken{ప్రమేయ సంకేతాలు}{!!{function}s} లేని \tetoken{సూత్రాలనే}{!!{formula}s} పరిశీలిస్తాం.

\begin{thm}[క్రెయిగ్ అంతర్వర్తన ప్రమేయం]\ollabel{thm:interpolation}
$!A$ యొక్క \tetoken{భాష}{!!{language}} $\Lang L(!A)$, $!B$ది $\Lang L(!B)$ అనుకుందాం. $\Log{G3c} \Proves !A \Sequent !B$ అయితే $\Lang L(!C) \subseteq \Lang
L(!A) \cap \Lang L(!B)$ \tetoken{భాషలో}{!!{language}} గల \tetoken{ఒక సూత్రం}{!!a{formula}}~$!C$ ఉంటుంది; దానికి $\Log{G3c} \Proves !A \Sequent !C$, $\Log{G3c} \Proves !C \Sequent !B$.
\end{thm}

~\Log{G3c}లోని \tetoken{నిరూపణల}{!!{proof}s} నిర్మాణాన్ని ఉపయోగించడానికి అంతర్వర్తన ప్రమేయం సాధారణీకరణాన్ని నిరూపిద్దాం. అందుకోసం సీక్వెంట్‌ల పూర్వాంగాలు, ఉత్తరాంగాలు రెండేసి భాగాలుగా విడిపోయినట్లు ఊహించండి. $\maeh{\Gamma_1}{\Gamma_2} \Sequent
\maeh{\Delta_1}{\Delta_2}$ అని రాస్తే $\Gamma_1$, $\Delta_1$ మొదటి భాగానికి, $\Gamma_2$, $\Delta_2$ రెండో భాగానికి చెందినవని అర్థం. అప్పుడు కింది లెమ్మా నుంచి అంతర్వర్తన ప్రమేయం వెంటనే వస్తుంది.

\begin{lem}[మహేరా లెమ్మా]\ollabel{lem:maehara} $\Log{G3c}
\Proves \maeh{\Gamma_1}{\Gamma_2} \Sequent \maeh{\Delta_1}{\Delta_2}$ అయితే $\Lang L(!C)
\subseteq \Lang L(\Gamma_1,
\Delta_1) \cap \Lang L(\Gamma_2, \Delta_2)$ \tetoken{భాషలో}{!!{language}} గల \tetoken{ఒక సూత్రం}{!!a{formula}}~$!C$ ఉంటుంది; దానికి $\Log{G3c}
\Proves \Gamma_1 \Sequent \Delta_1, !C$, $\Log{G3c} \Proves !C,
\Gamma_2 \Sequent \Delta_2$.
\end{lem}

\begin{proof}
  $\pi \Proves \maeh{\Gamma_1}{\Gamma_2} \Sequent
  \maeh{\Delta_1}{\Delta_2}$ అనుకుందాం. $n=\pheight{\pi}$పై ఆగమన పద్ధతిలో వాదనను నిరూపిద్దాం.

  $n=0$ అయితే $\maeh{\Gamma_1}{\Gamma_2} \Sequent
  \maeh{\Delta_1}{\Delta_2}$ ఒక స్వీకృతం; ఏదో అణు \tetoken{సూత్రం}{!!{formula}}~$!A$ అనేది $\Gamma_1, \Gamma_2$లోనూ $\Delta_1,
  \Delta_2$లోనూ ఉంటుంది. $!A$ ఎడమవైపు మొదటి లేదా రెండో భాగంలో, కుడివైపు మొదటి లేదా రెండో భాగంలో ఉండడాన్ని బట్టి నాలుగు సందర్భాలు.
  \begin{enumerate}
    \item $!A \in \Gamma_1$, $!A \in \Delta_1$. $\Gamma_1 \Sequent \Delta_1, !C$, $!C, \Gamma_2
    \Sequent \Delta_2$ రెండూ \tetoken{నిరూపణీయాలు}{!!{provable}} అయ్యే~$!C$ కావాలి. $!A$ ఎడమ, కుడి రెండింటిలోనూ ఉంది కనుక మొదటిది~$!C$ను ఎలా ఎంచుకున్నా స్వీకృతమే. $!C, \Gamma_2 \Sequent \Delta_2$ \tetoken{నిరూపణీయం}{!!{provable}} కావడానికి హామీ ఇచ్చే ఎంపిక $!C \ident \lfalse$.
    \item $!A \in \Gamma_1$, $!A \in \Delta_2$. అప్పుడు $\Gamma_1
    \Sequent \Delta_1, !A$, $!A, \Gamma_2 \Sequent \Delta_2$ రెండూ స్వీకృతాలు; అందువల్ల $!C \ident !A$ ఎంచుకోవచ్చు.
    \item $!A \in \Gamma_2$, $!A \in \Delta_1$. అప్పుడు $!A, \Gamma_1
    \Sequent \Delta_1$, $\Gamma_2 \Sequent \Delta_2, !A$ స్వీకృతాలు, కానీ వాటిలో~$!A$ తప్పు వైపున ఉంది. కాబట్టి $!A$ మధ్యవర్తి కాదు. అయితే $\RightR{\lnot}$, $\LeftR{\lnot}$లతో $\Gamma_1
    \Sequent \Delta_1, \lnot!A$, $\lnot !A, \Gamma_2 \Sequent
    \Delta_2$లను \tetoken{నిరూపించవచ్చు}{!!{prove}}.
    \item $!A \in \Gamma_2$, $!A \in \Delta_2$. సాధన.
  \end{enumerate}
  $\lfalse \in \Gamma_1$ లేదా $\lfalse \in \Gamma_2$ వల్ల $\maeh{\Gamma_1}{\Gamma_2} \Sequent
  \maeh{\Delta_1}{\Delta_2}$ స్వీకృతమయ్యే సందర్భాలను సాధనలుగా వదిలేస్తాం.

  $n>0$ అయితే $\pi$ ఒక తార్కిక నిగమనంతో ముగుస్తుంది. ఉపయోగించిన నియమం, ప్రధాన సూత్రం చెందిన భాగం బట్టి సందర్భాలను వేరుచేయాలి. ప్రతి సందర్భంలో $\pheight{\pi_1} < n$ గల అన్ని \tetoken{నిరూపణలకు}{!!{proof}s}~$\pi_1$, ~$\pi_1$ తుది సీక్వెంట్‌ను రెండు భాగాలుగా ఎలా విభజించినా, లెమ్మా వర్తిస్తుందనే ఆగమన పరికల్పనను ఉపయోగించవచ్చు. కొన్ని నిర్దిష్ట సందర్భాలను చూద్దాం.
  \begin{enumerate}
    \item $\pi$ చివర్లో ప్రధాన సూత్రం~$\lnot
    !A$ గల $\RightR{\lnot}$ ఉంది. $\lnot !A$ మొదటి భాగానిదా, రెండవ భాగానిదా అనే రెండు సందర్భాల ప్రకారం $\pi$ ఇలా ముగుస్తుంది:
    \[
    \AxiomC{}
    \RightLabel{$\pi_1$}
    \Deduce$\maeh{!A, \Gamma_1}{\Gamma_2} \fCenter
      \maeh{\Delta_1}{\Delta_2}$
    \RightLabel{\RightR{\lnot}}
    \UnaryInf$\maeh{\Gamma_1}{\Gamma_2} \fCenter
      \maeh{\Delta_1, \lnot !A}{\Delta_2}$
    \DisplayProof
    \]
    లేదా
    \[
    \AxiomC{}
    \RightLabel{$\pi_1$}
    \Deduce$\maeh{\Gamma_1}{!A, \Gamma_2} \fCenter
      \maeh{\Delta_1}{\Delta_2}$
    \RightLabel{\RightR{\lnot}}
    \UnaryInf$\maeh{\Gamma_1}{\Gamma_2} \fCenter
      \maeh{\Delta_1}{\Delta_2, \lnot !A}$
    \DisplayProof
    \]
    ప్రధాన \tetoken{సూత్రం}{!!{formula}}~$\lnot !A$ మొదటి భాగంలో ఉన్నప్పుడు, పూర్వాపేక్షలోని చురుకైన \tetoken{సూత్రం}{!!{formula}}~$!A$ను కూడా మొదటి భాగానికే కేటాయించామని గమనించండి. ఈ ఎంపిక మన చేతిలోనే ఉంది. చురుకైన సూత్రాన్ని మరొక భాగానికి కేటాయించినా ఆగమన పరికల్పన వర్తిస్తుంది; కానీ అప్పుడు మధ్యవర్తిని కనుగొనలేము (సాధన: ప్రయత్నించండి).

    మొదటి సందర్భంలో $\lnot !A$ మొదటి భాగానికి చెందుతుంది. ఆగమన పరికల్పన ప్రకారం~$\pi_1$ తుది సీక్వెంట్‌కు ఒక మధ్యవర్తి, అంటే ఈ రెండూ నెరవేరే \tetoken{ఒక సూత్రం}{!!a{formula}}~$!C_1$ ఉంది:
    \begin{align*}
    !A, \Gamma_1 & \Sequent \Delta_1, !C_1 &&\text{మరియు} &
    !C_1, \Gamma_2 & \Sequent \Delta_2
    \intertext{ఇవి నిరూపణీయాలు. మనకు కావలసింది ఈ క్రింది రెండూ \tetoken{నిరూపణీయాలు}{!!{provable}} అయ్యే \tetoken{ఒక సూత్రం}{!!a{formula}}~$!C$:}
    \Gamma_1 & \Sequent \Delta_1, \lnot !A, !C &&\text{మరియు} &
    !C, \Gamma_2 & \Sequent \Delta_2
    \end{align*}
    స్పష్టంగా $!C \ident !C_1$నే తీసుకోవచ్చు: కుడి సీక్వెంట్ \tetoken{నిరూపణీయం}{!!{provable}} అని ఆగమన పరికల్పన ఇప్పటికే చెబుతుంది; ఎడమ సీక్వెంట్ దాని నుంచి లభించిన సీక్వెంట్‌కు~$\RightR{\lnot}$ వర్తింపజేస్తే వస్తుంది. అలాగే ఆగమన పరికల్పన $\Lang L(!C_1) \subseteq \Lang L(!A, \Gamma_1,
    \Delta_1) \cap \Lang L(\Gamma_2, \Delta_2)$ అని చెబుతుంది. $\Lang L(\lnot
    !A) = \Lang L(!A)$, $\Lang L(!C) = \Lang L(!C_1)$ కనుక $\Lang L(!C_1) \subseteq \Lang L(\Gamma_1, \Delta_1, \lnot !A) \cap \Lang
    L(\Gamma_2, \Delta_2)$.


    రెండో సందర్భం కొంత భిన్నం: $\lnot !A$ రెండో భాగానికి చెందుతుంది. పూర్వాపేక్షలోని చురుకైన \tetoken{సూత్రాన్ని}{!!{formula}} కూడా రెండో భాగానికే కేటాయించి, ఆగమన పరికల్పనను వర్తింపజేస్తాం:
    \begin{align*}
    \Gamma_1 & \Sequent \Delta_1, !C_1 &&\text{మరియు} &
    !C_1, !A, \Gamma_2 & \Sequent \Delta_2
    \intertext{ఇవి నిరూపణీయాలు. వీటిలో రెండవదానికి మళ్లీ \RightR{\lnot} వర్తింపజేస్తే కింది \tetoken{నిరూపణలు}{!!{proof}s} వస్తాయి:}
    \Gamma_1 & \Sequent \Delta_1, !C_1 &&\text{మరియు} &
    !C_1, \Gamma_2 & \Sequent \Delta_2, \lnot !A
    \end{align*}
    $!C_1$ మధ్యవర్తి అయితే ఇవే మనకు కావలసిన \tetoken{నిరూపణీయ}{!!{provable}} సీక్వెంట్‌లు; మళ్లీ $!C \ident !C_1$ ఎంచుకుంటాం. పూర్వంలాగే $\Lang
    L(!C_1) \subseteq \Lang L(\Gamma_1, \Delta_1) \cap \Lang L(\Gamma_2,
    \Delta_2, \lnot !A)$.
    \item $\pi$ చివర్లో ప్రధాన \tetoken{సూత్రం}{!!{formula}}~$!A
    \lif !B$ గల $\RightR{\lif}$ ఉంది. $!A \lif !B$ మొదటి లేదా రెండో భాగానికి చెందడాన్ని బట్టి మళ్లీ రెండు సందర్భాలు. \tetoken{నిరూపణ}{!!{proof}}~$\pi$ ఇలా ముగుస్తుంది:
    \[
    \AxiomC{}
    \RightLabel{$\pi_1$}
    \Deduce$\maeh{!A, \Gamma_1}{\Gamma_2} \fCenter
      \maeh{\Delta_1, !B}{\Delta_2}$
    \RightLabel{\RightR{\lif}}
    \UnaryInf$\maeh{\Gamma_1}{\Gamma_2} \fCenter
      \maeh{\Delta_1, !A \lif !B}{\Delta_2}$
    \DisplayProof
    \]
    లేదా
    \[
    \AxiomC{}
    \RightLabel{$\pi_1$}
    \Deduce$\maeh{\Gamma_1}{!A, \Gamma_2} \fCenter
      \maeh{\Delta_1}{\Delta_2, !B}$
    \RightLabel{\RightR{\lif}}
    \UnaryInf$\maeh{\Gamma_1}{\Gamma_2} \fCenter
      \maeh{\Delta_1}{\Delta_2, !A \lif !B}$
    \DisplayProof
    \]
    మళ్లీ చురుకైన \tetoken{సూత్రాలను}{!!{formula}s} పూర్వాపేక్షలో ప్రధాన \tetoken{సూత్రం}{!!{formula}} ముగింపులో ఉన్న అదే భాగానికి కేటాయించాం. మొదటి సందర్భంలో ఆగమన పరికల్పన ద్వారా వీటికి \tetoken{నిరూపణలు}{!!{proof}s} వస్తాయి:
    \begin{align*}
    !A, \Gamma_1 & \Sequent \Delta_1, !B, !C_1 &&\text{మరియు} & !C_1, \Gamma_2 & \Sequent \Delta_2
    \intertext{ఎడమ సీక్వెంట్ \RightR{\lif}కు సరైన పూర్వాపేక్ష. కాబట్టి కింది సీక్వెంట్‌లు \tetoken{నిరూపణీయాలు}{!!{provable}}:}
    \Gamma_1 & \Sequent \Delta_1, !A \lif !B, !C_1 &&\text{మరియు} & !C_1, \Gamma_2 & \Sequent \Delta_2
    \end{align*}
    అందుచేత $!C_1$ అనేది~$\pi$ తుది సీక్వెంట్‌కు కూడా మధ్యవర్తి; మళ్లీ $!C \ident !C_1$ తీసుకోవచ్చు.
    
    మరొక సందర్భం కూడా అదే విధంగా సాగుతుంది.
    \item $\pi$ చివర్లో ప్రధాన \tetoken{సూత్రం}{!!{formula}}~$!A
    \lif !B$ గల $\LeftR{\lif}$ ఉంది. $!A \lif !B$ మొదటి భాగానికి చెందితే \tetoken{నిరూపణ}{!!{proof}}~$\pi$ ఇలా ముగుస్తుంది:
    \[
    \AxiomC{}
    \RightLabel{$\pi_1$}
    \Deduce$\maeh{\Gamma_1}{\Gamma_2} \fCenter
      \maeh{\Delta_1, !A}{\Delta_2}$
    \AxiomC{}
    \RightLabel{$\pi_2$}
    \Deduce$\maeh{!B, \Gamma_1}{\Gamma_2} \fCenter
      \maeh{\Delta_1}{\Delta_2}$
    \BinaryInf$\maeh{!A \lif !B, \Gamma_1}{\Gamma_2} \fCenter
      \maeh{\Delta_1}{\Delta_2}$
    \DisplayProof
    \]
    ఆగమన పరికల్పన ఎడమ పూర్వాపేక్ష మధ్యవర్తి $!C_1$కు రెండు, రెండో పూర్వాపేక్ష మధ్యవర్తి $!C_2$కు రెండు, మొత్తం నాలుగు \tetoken{నిరూపణీయ}{!!{provable}} సీక్వెంట్‌లను ఇస్తుంది:
    \begin{align*}
    \Gamma_1 & \Sequent \Delta_1, !A, !C_1 &&\text{(1)} &
    !C_1, \Gamma_2 & \Sequent \Delta_2 && \text{(2)}\\
    !B, \Gamma_1 & \Sequent \Delta_1, !C_2 &&\text{(3)} &
    !C_2, \Gamma_2 & \Sequent \Delta_2 && \text{(4)}
    \end{align*}
    (1), (3)లను బలహీనీకరించి, (1) కుడివైపు $!C_2$ను, (3) కుడివైపు $!C_1$ను చేరిస్తే, అవి~\LeftR{\lif}కు పూర్వాపేక్షలవుతాయి:
    \[
    \AxiomC{}
    \Deduce$\Gamma_1 \fCenter \Delta_1, !A, !C_1, !C_2$
    \AxiomC{}
    \Deduce$!B, \Gamma_1 \fCenter \Delta_1, !C_1, !C_2$
    \RightLabel{\LeftR{\lif}}
    \BinaryInf$!A \lif !B, \Gamma_1 \fCenter \Delta_1, !C_1, !C_2$
    \DisplayProof
    \]
    $\RightR{\lor}$ను వర్తింపజేస్తే ఈ సీక్వెంట్ వస్తుంది:
    \[
    !A \lif !B, \Gamma_1 \fCenter \Delta_1, !C_1 \lor !C_2
    \]
    కాబట్టి ఈ సందర్భంలో $!C_1 \lor !C_2$ మధ్యవర్తి కావచ్చని సూచిస్తుంది. నిజంగానే, మిగిలిన (2),~(4)ల నుంచి అవసరమైన మరొక సీక్వెంట్‌ను \tetoken{నిరూపించవచ్చు}{!!{prove}}:
    \[
    \AxiomC{}
    \Deduce$!C_1, \Gamma_2 \fCenter \Delta_2$
    \AxiomC{}
    \Deduce$!C_2, \Gamma_2 \fCenter \Delta_2$
    \RightLabel{\LeftR{\lor}}
    \BinaryInf$!C_1 \lor !C_2, \Gamma_2 \fCenter \Delta_2$
    \DisplayProof
    \]
    $!C_1 \lor !C_2$ సరైన భాషలో ఉందని తనిఖీ చేయాలి. ఆగమన పరికల్పన ప్రకారం:
    \begin{align*}
    \Lang L(!C_1) & \subseteq
    \Lang L(\Gamma_1, \Delta_1, !A) \cap \Lang L(\Gamma_2, \Delta_2) &\text{మరియు}\\
    \Lang L(!C_2) & \subseteq
    \Lang L(!B, \Gamma_1, \Delta_1) \cap \Lang L(\Gamma_2, \Delta_2)
    \intertext{ఎందుకంటే}
    \Lang L(!A) & \subseteq \Lang L(!A \lif !B) &\text{మరియు}\\
    \Lang L(!B) & \subseteq \Lang L(!A \lif !B)
    \intertext{రెండూ లభిస్తాయి:}
    \Lang L(!C_1) & \subseteq
    \Lang L(\Gamma_1, \Delta_1, !A \lif !B) \cap \Lang L(\Gamma_2, \Delta_2) &\text{మరియు}\\
    \Lang L(!C_2) & \subseteq
    \Lang L(\Gamma_1, \Delta_1, !A \lif !B) \cap \Lang L(\Gamma_2, \Delta_2)
    \intertext{అందువల్ల $\Lang L (!C_1 \lor !C_2) = {}$}
    \Lang L(!C_1) \cup \Lang L(!C_2) & \subseteq
    \Lang L(\Gamma_1, \Delta_1, !A \lif !B) \cap \Lang L(\Gamma_2, \Delta_2),
    \end{align*}
    ఇదే కావలసింది.

    $!A \lif !B$ రెండో భాగానికి చెందిన సందర్భం భిన్నమైనా అదే తరహాలో సాగుతుంది. ఆగమన పరికల్పన మళ్లీ నాలుగు \tetoken{నిరూపణీయ}{!!{provable}} సీక్వెంట్‌లను ఇస్తుంది:
    \begin{align*}
    \Gamma_1 & \Sequent \Delta_1, !C_1 &&\text{(1)} &
    !C_1, \Gamma_2 & \Sequent \Delta_2, !A && \text{(2)}\\
    \Gamma_1 & \Sequent \Delta_1, !C_2 &&\text{(3)} &
    !C_2, !B, \Gamma_2 & \Sequent \Delta_2 && \text{(4)}
    \end{align*}
    ఇప్పుడు (2), (4) సీక్వెంట్‌లు—కొన్ని బలహీనీకరించిన \tetoken{సూత్రాలతో}{!!{formula}s}—\LeftR{\lif}కు పూర్వాపేక్షలవుతాయి:
    \[
    \AxiomC{}
    \Deduce$!C_1, !C_2, \Gamma_2 \fCenter \Delta_2, !A$
    \AxiomC{}
    \Deduce$!B, !C_1, !C_2, \Gamma_2 \fCenter \Delta_2$
    \RightLabel{\LeftR{\lif}}
    \BinaryInf$!A \lif !B, !C_1, !C_2, \Gamma_2 \fCenter \Delta_2$
    \DisplayProof
    \]
    దీన్నీ $!C_1$, $!C_2$ల నుంచి నిర్మించిన ఒకే \tetoken{సూత్రం}{!!{formula}} గల సీక్వెంట్‌గా మార్చాలి; కానీ ఇప్పుడు అది పూర్వాంగంలో కావాలి (రెండో భాగం గురించి మాట్లాడుతున్నాం కనుక). $\LeftR{\land}$ వర్తింపజేస్తే సరిపోతుంది; $!C \ident
    (!C_1 \land !C_2)$ను మధ్యవర్తిగా తీసుకుందాం. మిగిలిన (1), (3)లను ఉపయోగించి మరొక భాగానికి అనుగుణ సీక్వెంట్‌ను \tetoken{నిరూపించవచ్చు}{!!{prove}}:
    \[
    \AxiomC{}
    \Deduce$\Gamma_1 \fCenter \Delta_1, !C_1$
    \AxiomC{}
    \Deduce$\Gamma_1 \fCenter \Delta_1, !C_2$
    \RightLabel{\RightR{\land}}
    \BinaryInf$\Gamma_1 \fCenter \Delta_1, !C_1 \land !C_2$
    \DisplayProof
    \]
    పూర్వంలాగే $\Lang L(!C_1 \land !C_2) \subseteq \Lang
    L(\Gamma_1, \Delta_1) \cap \Lang L(\Gamma_2, \Delta_2, !A \lif
    !B)$.

    \item $\pi$ చివర్లో ప్రధాన \tetoken{సూత్రం}{!!{formula}}~$\lforall[x][!A(x)]$ గల $\RightR{\lforall}$ ఉంది:
    \[
    \AxiomC{}
    \RightLabel{$\pi_1$}
    \Deduce$\maeh{\Gamma_1}{\Gamma_2} \fCenter
    \maeh{\Delta_1, !A(c)}{\Delta_2}$
    \RightLabel{\RightR{\lforall}}
    \UnaryInf$\maeh{\Gamma_1}{\Gamma_2} \fCenter
      \maeh{\Delta_1, \lforall[x][!A(x)]}{\Delta_2}$
    \DisplayProof
    \]
    లేదా
    \[
    \AxiomC{}
    \RightLabel{$\pi_1$}
    \Deduce$\maeh{\Gamma_1}{\Gamma_2} \fCenter
      \maeh{\Delta_1}{\Delta_2, !A(c)}$
    \RightLabel{\RightR{\lforall}}
    \UnaryInf$\maeh{\Gamma_1}{\Gamma_2} \fCenter
      \maeh{\Delta_1}{\Delta_2, \lforall[x][!A(x)]}$
    \DisplayProof
    \]
    మొదటి సందర్భంలో ఆగమన పరికల్పన $\Gamma_1 \Sequent \Delta_1, !A(c), !C_1$, $!C_1, \Gamma_2
    \Sequent \Delta_2$లకు \tetoken{నిరూపణలు}{!!{proof}s} ఇస్తుంది. మొదటిదానికి $\RightR{\lforall}$ వర్తింపజేసి $\Gamma_1 \Sequent \Delta_1, \lforall[x][!A(x)], !C_1$ పొందవచ్చు. (ఈజెన్‌చరరాశి షరతు నెరవేరుతుంది;~$\pi$ చివరి \RightR{\lforall} నియమంలోనే అది నెరవేరింది.) అందువల్ల $\Lang
    L(!C_1) \subseteq \Lang L(\Gamma_1, \Delta_1, \lforall[x][!A(x)])
    \cap \Lang L(\Gamma_2, \Delta_2)$ షరతు నెరవేరితే $!C_1$ మధ్యవర్తి. ఆగమన పరికల్పన ప్రకారం $\Lang L(!C_1) \subseteq \Lang L(\Gamma_1,
    \Delta_1, !A(c)) \cap \Lang L(\Gamma_2, \Delta_2)$. కాబట్టి~$c$ గురించి ప్రశ్న: $c$ అనేది~$!C_1$లో కనిపించగలదా? కాదు. $c$ అనేది $!C_1$లో కనిపిస్తే $c \in \Lang L(\Gamma_1, \Delta_1,
    !A(c))$ (ఇది నిజమే)తో పాటు $c \in \Lang L(\Gamma_2,
    \Delta_2)$ కూడా ఉండాలి. అప్పుడు~$\pi$లోని $\RightR{\lforall}$ నిగమనం ముగింపులో $c$ కనిపించి, ఈజెన్‌చరరాశి షరతు ఉల్లంఘిస్తుంది.
    
    మరొక సందర్భం కూడా సదృశమే.

    \item $\pi$ చివర్లో ప్రధాన \tetoken{సూత్రం}{!!{formula}}~$\lexists[x][!A(x)]$ గల $\RightR{\lexists}$ ఉంది. మళ్లీ రెండు సందర్భాలున్నాయి. $\lexists[x][!A(x)]$ మొదటి భాగానికి చెందిన సందర్భాన్ని పరిశీలించి, రెండోదాన్ని సాధనగా వదిలేస్తాం.

    మొదటి సందర్భంలో $\pi$ ఇలా ముగుస్తుంది:
    \[
    \AxiomC{}
    \RightLabel{$\pi_1$}
    \Deduce$\maeh{\Gamma_1}{\Gamma_2} \fCenter
      \maeh{\Delta_1, \lexists[x][!A(x)], !A(t)}{\Delta_2}$
    \RightLabel{\RightR{\lexists}}
    \UnaryInf$\maeh{\Gamma_1}{\Gamma_2} \fCenter
      \maeh{\Delta_1, \lexists[x][!A(x)]}{\Delta_2}$
    \DisplayProof
    \]
    ఆగమన పరికల్పన $\Gamma_1 \Sequent
    \Delta_1, \lexists[x][!A(x)], !A(t), !C_1$, $!C_1, \Gamma_2
    \Sequent \Delta_2$లకు \tetoken{నిరూపణలు}{!!{proof}s} ఇస్తుంది. మొదటిదానికి $\RightR{\lexists}$ వర్తింపజేసి $\Gamma_1 \Sequent \Delta_1, \lexists[x][!A(x)], !C_1$ పొందవచ్చు. కాబట్టి $\Lang L(!C_1) \subseteq \Lang L(\Gamma_1, \Delta_1,
    \lexists[x][!A(x)]) \cap \Lang L(\Gamma_2, \Delta_2)$ నెరవేరితే $!C_1$ మళ్లీ మధ్యవర్తి. ఇక్కడ ఆగమన పరికల్పన $\Lang L(!C_1)
    \subseteq \Lang L(\Gamma_1, \Delta_1, \lexists[x][!A(x)], !A(t))
    \cap \Lang L(\Gamma_2, \Delta_2)$ను ఇస్తుంది. \tetoken{ప్రమేయ సంకేతాలు}{!!{function}s} లేవని అనుకున్నాం గుర్తుంచుకోండి. అందువల్ల పదం~$t$ చరరాశి లేదా \tetoken{ఒక స్థిరసంకేతం}{!!a{constant}} కావచ్చు. చరరాశి భాషా సంకేతాల సమితిలో చేరదు. పరిశీలించాల్సిన స్థిరసంకేత సందర్భంలో $t \ident b$. \sourcecorrection{OLTEPTCUTITP-005}{ప్రమేయ సంకేతాలు లేకపోయినా పదం చరరాశి కావచ్చు; మూలం స్థిరసంకేత సందర్భమే సాధ్యమని పేర్కొంది. చరరాశి అయితే అదనపు భాషా చిహ్నం రాదని వేరుచేశాం.} $b \notin \Lang L(!C_1)$ అయినా, లేదా $b \in \Lang L(\Gamma_1, \Delta_1, \lexists[x][!A(x)]) \cap
    \Lang L(\Gamma_2, \Delta_2)$ అయినా సమస్య లేదు: $!C_1$నే మధ్యవర్తిగా తీసుకోవచ్చు. కానీ $b \in \Lang
    L(!C_1)$, $b \notin \Lang L(\Gamma_1, \Delta_1,
    \lexists[x][!A(x)]) \cap \Lang L(\Gamma_2, \Delta_2)$ అయితే? మొదట~$!C_1$లోని $b$ సంభవాలన్నిటినీ \tetoken{ఒక చరరాశి}{!!a{variable}}~$y$తో భర్తీ చేస్తే $!C_1 \ident \Subst{!C_1'}{b}{y}$ అని రాయవచ్చు; ఇక్కడ~$!C_1'$లో $b$ ఉండదు. ఇంకా $b \notin
    \Lang L(\Gamma_1, \Delta_1, \lexists[x][!A(x)])$ లేదా $b \notin
    \Lang L(\Gamma_2, \Delta_2)$. వరుసగా $!C \ident
    \lforall[y][!C_1'(y)]$ లేదా $!C \ident \lexists[y][!C_1'(y)]$ తీసుకోవచ్చు. మొదటి సందర్భంలో:
    \[
    \AxiomC{}
    \Deduce$\Gamma_1 \fCenter
      \Delta_1, \lexists[x][!A(x)], !A(b), !C_1'(b)$
    \RightLabel{\RightR{\lexists}}
    \UnaryInf$\Gamma_1 \fCenter \Delta_1, \lexists[x][!A(x)], !C_1'(b)$
    \RightLabel{\RightR{\lforall}}
    \UnaryInf$\Gamma_1 \fCenter
      \Delta_1, \lexists[x][!A(x)], \lforall[y][!C_1'(y)]$
    \DisplayProof
    \]
    మరియు
    \[
    \AxiomC{}
    \Deduce$!C_1'(b), \lforall[y][!C_1'(y)], \Gamma_2 \fCenter \Delta_2$
    \RightLabel{\LeftR{\lforall}}
    \UnaryInf$\lforall[y][!C_1'(y)], \Gamma_2 \fCenter \Delta_2$
    \DisplayProof
    \]
    ఎగువ సీక్వెంట్‌ల \tetoken{నిరూపణలు}{!!{proof}s} ఆగమన పరికల్పన నుంచి వస్తాయి (రెండోదానిలో పూర్వాంగానికి $\lforall[y][!C_1'(y)]$ చేర్చే బలహీనీకరణ అనుమతిత్వాన్ని వాడుతాం). మొదటి \tetoken{నిరూపణలో}{!!{proof}} \RightR{\lforall} ఈజెన్‌చరరాశి షరతు నెరవేరుతుంది, ఎందుకంటే $b \notin \Lang L(\Gamma_1, \Delta_1,
    \lexists[x][!A(x)])$; అంతేకాక~$!C_1'$లో $b$ ఉండకుండా నిర్వచించాం.
    
    రెండో సందర్భంలో (ఆగమన పరికల్పనతో పాటు~$\lexists[y][!C_1'(y)]$ ద్వారా బలహీనీకరణ అనుమతిత్వాన్ని ఉపయోగించి) ఇవి లభిస్తాయి:
    \[
    \AxiomC{}
    \Deduce$\Gamma_1 \fCenter
      \Delta_1, \lexists[x][!A(x)], !A(b), \lexists[y][!C_1'(y)], !C_1'(b)$
    \RightLabel{\RightR{\lexists}}
    \UnaryInf$\Gamma_1 \fCenter
      \Delta_1, \lexists[x][!A(x)], \lexists[y][!C_1'(y)], !C_1'(b)$
    \RightLabel{\RightR{\lexists}}
    \UnaryInf$\Gamma_1 \fCenter
      \Delta_1, \lexists[x][!A(x)], \lexists[y][!C_1'(y)]$
    \DisplayProof
    \]
    మరియు
    \[
    \AxiomC{}
    \Deduce$!C_1'(b), \Gamma_2 \fCenter \Delta_2$
    \RightLabel{\LeftR{\lexists}}
    \UnaryInf$\lexists[y][!C_1'(y)], \Gamma_2 \fCenter \Delta_2$
    \DisplayProof
    \]
    రెండో \tetoken{నిరూపణలో}{!!{proof}} \LeftR{\lexists} ఈజెన్‌చరరాశి షరతు నెరవేరుతుంది, ఎందుకంటే $b \notin \Lang L(\Gamma_2,
    \Delta_2)$; అలాగే~$!C_1'$లో $b$ ఉండకుండా నిర్వచించాం.
  \end{enumerate}
  మిగిలిన సందర్భాలు సదృశం; వాటిని సాధనలుగా వదిలేస్తాం.
\end{proof}

\begin{prob}
$\lnand$ అనే సంయోజకానికి ఈ నియమాలు ఉన్నాయనుకుందాం:
\[
\Axiom$\Gamma \fCenter \Delta, !A$
\Axiom$\Gamma \fCenter \Delta, !B$
\RightLabel{\LeftR{\lnand}}
\BinaryInf$!A \lnand !B, \Gamma \fCenter \Delta$
\DisplayProof
\quad
\Axiom$!A, !B, \Gamma \fCenter \Delta$
\RightLabel{\RightR{\lnand}}
\UnaryInf$\Gamma \fCenter \Delta, !A \lnand !B$
\DisplayProof
\]
$\pi$ ఈ నియమాల్లో ఒకదానితో ముగిసే \olref{lem:maehara} నిరూపణ సందర్భాలను పూర్తి చేసి, మహేరా లెమ్మా ఇంకా వర్తిస్తుందని చూపండి.
\end{prob}

మహేరా లెమ్మా స్థాపించాం కనుక దాన్ని క్రెయిగ్ అంతర్వర్తన ప్రమేయాన్ని నిరూపించడానికి ఉపయోగించవచ్చు.

\begin{proof}[\olref{thm:interpolation} నిరూపణ]
  $!A \Sequent !B$ \tetoken{నిరూపణీయం}{!!{provable}} అనుకుందాం. $\maeh{!A}{\ } \Sequent \maeh{\ }{!B}$కు \olref{lem:maehara} వర్తింపజేసి మధ్యవర్తి~$!C$ను పొందండి. అప్పుడు $\Proves !A \Sequent !C$, $\Proves !C
  \Sequent !B$.
\end{proof}

$\Gamma$ అనేది $n$-స్థాన {విధేయం}~$R$ కలిగిన \tetoken{ఒక భాష}{!!a{language}}~$\Lang L$లోని ఒక సిద్ధాంతం (\tetoken{వాక్యాల}{!!{sentence}s} సమితి) అనుకుందాం. $\Gamma$~$R$ను \emph{పరోక్షంగా నిర్వచిస్తుంది} అంటే:
\begin{align*}
\Gamma, \Gamma'
& \Entails \lforall[x_1][\dots\lforall[x_n][(R(x_1,\dots,x_n) \liff
R'(x_1, \dots, x_n))]],
\intertext{ఇక్కడ $\Gamma'$ అనేది $\Gamma$లోని~$R$ అన్ని సంభవాలను $R' \notin \Lang L$తో భర్తీ చేస్తే వచ్చేది. $\Gamma$~$R$ను \emph{ప్రత్యక్షంగా నిర్వచిస్తుంది} అంటే,~$R$ లేని ఏదో~$!A(x_1, \dots, x_n)$ ఉండి}
\Gamma &\Entails
\lforall[x_1][\dots\lforall[x_n][(R(x_1,\dots,x_n) \liff !A(x_1,
\dots, x_n))]].
\end{align*}
$\Gamma$~$R$ను ప్రత్యక్షంగా నిర్వచిస్తే~$R$ను పరోక్షంగానూ నిర్వచిస్తుందని స్పష్టం. దీనికి విపర్యయం అంత స్పష్టంగా కనిపించకపోయినా నిజమే:

\begin{lem}[బెత్ నిర్వచనీయత ప్రమేయం]
  $\Gamma$~$R$ను పరోక్షంగా నిర్వచిస్తే ప్రత్యక్షంగానూ నిర్వచిస్తుంది.
\end{lem}

\begin{proof}
  $\Gamma$~$R$ను పరోక్షంగా నిర్వచిస్తుందనుకుందాం. పైన చెప్పినట్లే $R'$, $\Gamma'$ తీసుకుని $\Lang L' = \Lang L \setminus \{R\}
  \cup \{R'\}$ అనుకుందాం. \sourcecorrection{OLTEPTCUTITP-006}{మూలం $\Lang L(\Gamma')=\Lang L'$ అని చెప్పింది; కానీ నిర్వచించిన $\Lang L(\Gamma')$లో నిజంగా కనిపించే సంకేతాలే ఉంటాయి. సిద్ధాంతం అన్ని భాషా సంకేతాలను ఉపయోగించాలనే షరతు లేదు; ఇక్కడ అవసరమైనది $\Gamma'$ లక్ష్య భాషలో ఉండడమే.} పరికల్పన ప్రకారం:
  \begin{align*}
    \Gamma, \Gamma' &
    \Sequent \lforall[x_1][\dots\lforall[x_n][(R(x_1,\dots,x_n)
    \liff R'(x_1, \dots, x_n))]]
  \intertext{~$\Gamma$లో కనిపించని \tetoken{స్థిరసంకేతాలు}{!!{constant}s} $c_1$, \dots,~$c_n$ తీసుకోండి. విలోమయోగ్యత లెమ్మా ప్రకారం:}
    \Gamma, \Gamma', R(c_1,\dots,c_n) & \Sequent  R'(c_1, \dots, c_n)
  \intertext{దీనికి మహేరా లెమ్మా వర్తింపజేయండి:}
    \maeh{\Gamma, R(c_1,\dots,c_n)}{\Gamma'} & \Sequent
      \maeh{\ }{R'(c_1, \dots, c_n)}
  \intertext{అప్పుడు $!A \ident !A(x_1, \dots, x_n)$ దొరుకుతుంది; దానికి:}
    \Gamma, R(c_1,\dots,c_n) & \Sequent !A(c_1, \dots, c_n) \\
    !A(c_1, \dots, c_n), \Gamma' & \Sequent R'(c_1, \dots, c_n)
  \intertext{వీటికి \tetoken{నిరూపణలు}{!!{proof}s} ఉన్నాయి. మహేరా భాషా పరిమితి, తాజా స్థిరసంకేతాలను చరరాశులుగా మార్చడం వల్ల $\Lang L(!A) \subseteq \Lang L(\Gamma) \setminus \{R\}$; కాబట్టి~$!A$లో $R$ లేదు. రెండో సీక్వెంట్ \tetoken{నిరూపణ}{!!{proof}}లో ప్రతిచోటా $R'$ స్థానంలో~$R$ను పెట్టి దీనికి \tetoken{ఒక నిరూపణ}{!!a{proof}} పొందండి:}
    !A(c_1, \dots, c_n), \Gamma & \Sequent R(c_1, \dots, c_n)
  \end{align*}
  \sourcecorrection{OLTEPTCUTITP-001}{మూలంలోని $\Lang L(!A)=\Lang L_1\cap\Lang L_2=\Lang L(\Gamma)\setminus\{R\}$ సమానతలో $L_1,L_2$ నిర్వచించబడలేదు; మహేరా లెమ్మా సమితి ఉపసమితి సంబంధాన్నే ఇస్తుంది. తాజా స్థిరసంకేతాలను చరాలుగా మార్చాక సరిపడే ఉపసమితి సంబంధాన్ని ఇక్కడ పేర్కొన్నాం.}
  \RightR{\lif}, \RightR{\lforall}లను వర్తింపజేసి దీనికి \tetoken{ఒక నిరూపణ}{!!a{proof}} పొందండి:
  \[
    \Gamma \Sequent \lforall[x_1][\dots\lforall[x_n][(R(x_1,\dots,x_n) \liff !A(x_1, \dots, x_n))]].
  \]
  ఇది $!A(x_1, \dots, x_n)$ అనేది~$\Gamma$లో $R$ను ప్రత్యక్షంగా నిర్వచిస్తుందని చూపుతుంది.
  \sourcecorrection{OLTEPTCUTITP-002}{చివరి వాక్యంలో మూలం విధేయ సంకేతం $R$ను అనవసరమైన ఆశ్చర్యార్థ చిహ్నంతో $!R$గా రాసింది; నిర్వచనంలోని సంకేతం $R$గా సరిచేశాం.}
\end{proof}

అంతర్వర్తనకు సమానమైన (దానిలాగే మహేరా లెమ్మాతో నిరూపించగల) మూడో ఫలితం ఇది: రెండు అవైరుధ్య సిద్ధాంతాలు~$\Gamma_1$, $\Gamma_2$ ఉమ్మడి భాష $\Lang L(\Gamma) = \Lang L(\Gamma_1)
\cap \Lang L(\Gamma_2)$లోని సంపూర్ణ సిద్ధాంతం $\Gamma$ను కలిగి ఉంటే, $\Gamma_1 \cup \Gamma_2$ అవైరుధ్యం.
\sourcecorrection{OLTEPTCUTITP-003}{మూలం ఇక్కడ ఉమ్మడి భాషతో సమానతకు బదులు ఉపసమితి షరతు మాత్రమే ఇచ్చింది. కానీ వెంట వచ్చే వాదన ఉమ్మడి భాషలోని ప్రతి వాక్యానికి $\Gamma$ సంపూర్ణతను వర్తింపజేస్తుంది; దానికి పూర్తి ఉమ్మడి భాష కావాలి.}

సంపూర్ణ సిద్ధాంతం అంటే దాని భాషలోని ప్రతి \tetoken{వాక్యం}{!!{sentence}}~$!A$కు $\Gamma \Proves !A$ లేదా $\Gamma \Proves \lnot !A$ ఉండే సిద్ధాంతం అని గుర్తుంచుకోండి. $\Gamma_1$,~$\Gamma_2$లు~$\Gamma$ యొక్క అవైరుధ్య విస్తరణలు కనుక $\Gamma$ అవైరుధ్యమే. ఇప్పుడు ఉమ్మడి భాష $\Lang L(\Gamma_1) \cap \Lang
L(\Gamma_2)$లో \tetoken{ఒక వాక్యం}{!!a{sentence}}~$!A$ తీసుకుని $\Gamma_1 \Entails !A$, $\Gamma_2
\Entails \lnot !A$ రెండూ నెరవేరే అలాంటి~$!A$ ఉందని విరుద్ధార్థం అనుకుందాం. $\Gamma$ సంపూర్ణత వల్ల $\Gamma \Entails !A$ లేదా $\Gamma \Entails \lnot !A$. రెండోది అయితే $\Gamma \subseteq \Gamma_1$ కనుక $\Gamma_1 \Entails \lnot !A$ కూడా వచ్చి $\Gamma_1$ అవైరుధ్యానికి విరుద్ధం. అందువల్ల $\Gamma \Entails !A$. $\Gamma
\subseteq \Gamma_2$ కనుక $\Gamma_2 \Entails !A$; ఇది $\Gamma_2 \Entails \lnot !A$తో కలిపి $\Gamma_2$ అవైరుధ్యానికి విరుద్ధం.
\sourcecorrection{OLTEPTCUTITP-004}{మూల వాదన $\Gamma_1$ నుంచి $!A$ అనుగమిస్తే $\Gamma$ నుంచి కూడా అనుగమిస్తుందని నేరుగా చెప్పింది; ఈ తిరుగు అనుగమనం సాధారణంగా చెల్లదు. సంపూర్ణతలోని రెండు ప్రత్యామ్నాయాలను పరీక్షించి, చివర్లో $\Gamma_2$ అవైరుధ్యానికి విరుద్ధతను చూపే క్రమాన్ని సరిచేశాం.}

ఉమ్మడి \tetoken{భాష}{!!{language}}~$\Lang L(\Gamma_1) \cap \Lang L(\Gamma_2)$లో $\Gamma_1
\Entails !A$, $\Gamma_2 \Entails \lnot !A$ అయ్యే \tetoken{వాక్యం}{!!a{sentence}}~$!A$ ఏదీ లేకపోవడమే కింది నిరూపణకు అవసరం. దాన్ని మహేరా లెమ్మాతో నిరూపణ-సిద్ధాంత పద్ధతిలో ఇస్తాం.

\begin{thm}[రాబిన్సన్ సంయుక్త అవైరుధ్య ప్రమేయం]
$\Gamma_1$, $\Gamma_2$ అవైరుధ్య సిద్ధాంతాలనుకుందాం. \tetoken{భాష}{!!{language}} $\Lang L(\Gamma_1) \cap \Lang L(\Gamma_2)$లోని ఏ~$!A$కూ $\Gamma_1 \Entails !A$, $\Gamma_2 \Entails \lnot !A$ రెండూ నిజం కావనుకుందాం. అప్పుడు $\Gamma_1 \cup \Gamma_2$ అవైరుధ్యం.
\end{thm}

\begin{proof}
  విరుద్ధంగా $\Gamma_1 \cup \Gamma_2$ వైరుధ్యమని అనుకుందాం. అప్పుడు $\Gamma_1, \Gamma_2 \Sequent \ $ సీక్వెంట్‌కు \tetoken{ఒక నిరూపణ}{!!a{proof}} ఉంటుంది. $\maeh{\Gamma_1}{\Gamma_2}
  \Sequent \maeh{\ }{\ }$కు మహేరా లెమ్మా వర్తింపజేస్తే ఉమ్మడి \tetoken{భాష}{!!{language}}~$\Lang L(\Gamma_1) \cap \Lang L(\Gamma_2)$లో \tetoken{ఒక వాక్యం}{!!a{sentence}}~$!A$ దొరుకుతుంది; దానికి $\Proves \Gamma_1 \Sequent !A$, $\Proves !A, \Gamma_2 \Sequent
  \quad$. రెండవదాని నుంచి $\Proves \Gamma_2 \Sequent \lnot !A$ వస్తుంది. కానీ అలాంటి \tetoken{వాక్యం}{!!{sentence}} ఉండదని అనుకున్నాం.
\end{proof}

పైన $\Gamma$, $\Gamma_1$, $\Gamma_2$ సిద్ధాంతాలు పరిమితమైనవని మౌనంగా అనుకున్నాం. సంహతత్వం వల్ల ఫలితాలు అపరిమిత సిద్ధాంతాలకూ వర్తిస్తాయి.

\end{document}

